\scanpage{416}
\mspage{330}
% archive sheet number: 415

\deleted{

% 1–2
2) \inserted{$\forall\,\alpha\in I_{\alpha_0}$,} $I$ est \unsure{grd|grand} devant \deleted{\unsure{tout cardinal} $\unsure{\pi'<\pi}$ \illegible{2}} \inserted{$\operatorname{Card}\alpha$, pour tt $\alpha\in I_{\alpha_0}$ ;} \deleted{(\unsure{ou} \illegible{3})}

% 3–4
1) \inserted{Le foncteur} $\varphi$ est cofinal. [NB On sait que \deleted{\unsure{cel}} \inserted{1) et 2)} \unsure{impliquent} que $J$ est grd devant \deleted{\unsure{tout cardinal} $\unsure{\pi'<\pi}$} \inserted{$\operatorname{Card}\alpha$, pour tt $\alpha\in I_{\alpha_0}$ ;} \inserted{(\unsure{resp.} \unsure{devant} $\unsure{\pi}$)} ]

% 5–8
2) $I$ est stable par $\varinjlim$ de type $I_{\alpha\deleted{\unsure{\pi}}}$ \deleted{(où $\alpha_\pi$ est le plus petit ordinal $\alpha$ \deleted{\unsure{tel}} que $\operatorname{card}\alpha=\pi$), i.e. par $\varinjlim$ indexées par des ens. \unsure{préord.} bien ordonnés} \deleted{\unsure{de cardinal} $\pi$}. et le foncteur $\varphi$ commute à ces limites.

% 9–13
3) Le foncteur $\psi$ est l'inclusion d'une sous-catégorie str. pleine $J'$ de $J$ ; \deleted{et $J'$ est} De plus $J'$ est stable par $\varinjlim$ du type \deleted{\unsure{envisagé}} $I_{\alpha_\pi}$, et \unsure{le|les} \unsure{foncteur|foncteurs} $\unsure{\psi|f}$ \unsure{commute|commutent} \unsure{aux} $\varinjlim$.

% 14–16
4) $J'$ est cofinale dans $J$. [Donc il est grand devant tt cardinal $\pi'<\pi$, \deleted{$J$ l'étant \unsure{(?)}} (resp. devant $\pi$) \unsure{$J$ l'étant.|$J$ l'étant aussi.}]

% 17–21
Sous ces conditions, $I'$ est cofinale dans $I$. [Et il s'ensuit aussitôt que $I'$, tout comme $I,J,J'$, est grd devant tt cardinal $\pi'<\pi$ (resp. devant $\pi$), et on voit aussitôt que le foncteur $I'\unsure{-|\to}J'$ est également cofinal.]
}
